\section*{Chapter \ref{ch:m1:getting-access-equities} --- Getting Access: Equity Markets}

\begin{solution}{exo:m1:getting-access-equities:1}
$9/12\,000 = 0.075\%$ of consolidated volume: above 0.06\%, below 0.20\%: the
first tier, \$0.0020 a share. \$18\,000 a day; \$378\,000 for the month.
\end{solution}

\begin{solution}{exo:m1:getting-access-equities:2}
$23.9 \times 10^6 \times 0.0020 = \$47\,800$; $24.0 \times 10^6 \times 0.0023 =
\$55\,200$. The last hundred thousand shares are worth \$7\,400 a day, 7.4 cents
a share against a rebate of 0.23 cent, and \$155\,400 over the month.
\end{solution}

\begin{solution}{exo:m1:getting-access-equities:3}
Passive: $-20 + 2 + 1 = -17$ cents per 100 shares (income). Aggressive: $30 + 2
+ 1 = +33$. Round trip: a net cost of 16 cents per 100 shares, which a
strategy entering passively and exiting aggressively must earn before it
earns anything.
\end{solution}

\begin{solution}{exo:m1:getting-access-equities:4}
The larger of the dealer minimum, \$100\,000, and the market-maker amount, $300
\times 2\,500 = \$750\,000$ (under the \$1 million cap): \$750\,000. Maximum
aggregate indebtedness: $15 \times 2$ million $= \$30$ million. Neither binds in
practice: net capital is computed after haircuts on the firm's positions, so
a book of a few tens of millions consumes the \$2 million long before
indebtedness does, and the \omterm{def:m1:getting-access-equities:brokers}{clearing broker}'s house margin is usually
stricter still.
\end{solution}

\begin{solution}{exo:m1:getting-access-equities:5}
$10\% \times 390 = 39$ minutes. With ``a bid \emph{or} an offer at the inside'',
6\% and 7\% that never coincide make 13\% of the day with one side at the
inside: qualified under that reading; if they always coincide, 7\%: not
qualified. If the exchange instead averages the bid and offer percentages,
the first firm has 6.5\% and fails. The measurement method is in the rule and
must be read, since it decides how the quoting engine is tuned.
\end{solution}

\begin{solution}{exo:m1:getting-access-equities:6}
Queue position: few participants rest orders where they must pay, so the
queue is short and an order there is filled before the long queues on
rebate-paying venues, when the fill is most valuable. Flow: marketable
orders from fee-sensitive routers go to the \omterm{def:m1:us-equity-market-structure:makertaker}{inverted venue} first, because it
pays them, so the resting order meets a different, often less informed,
counterparty. Both can be worth more than 30 cents per 100 shares.
\end{solution}

\begin{solution}{exo:m1:getting-access-equities:7}
$24 \times (27 - 23)/(27 - 20) = 13.7$ million shares a day at a loss of 27;
$24 \times (25 - 23)/(25 - 20) = 9.6$ million at 25. Just above each, padding
gains a few dollars a day; just below, it loses them.
\end{solution}

\begin{solution}{exo:m1:getting-access-equities:8}
(i) The 31-cent rate is the top tier, for members adding 1\% of consolidated
volume, 120 million shares a day at 12 billion; at 2 million the firm is at
the standard 16. (ii) Clearing and regulatory costs, about 3, are missing.
Recomputed: $10 + 16 - 3 = 23$ cents per 100 shares, \$96\,600 a month, not
\$172\,200. One should add that the 10 cents ``before fees'' was presumably
measured on fills that a rebate-paying venue's long queue would not give.
\end{solution}

\begin{solution}{pb:m1:getting-access-equities:1}
\textbf{1.} $20/12\,000 = 0.167\%$: first tier, \$0.0020: \$840\,000 a month.
\textbf{2.} $0.20\% \times 12$ billion $= 24$ million: 4 million more a day.
\textbf{3.} $24 \times 0.0023 - 4 \times 0.0027 - 20 \times 0.0020$ (in millions of
dollars) $= 55\,200 - 10\,800 - 40\,000 = +\$4\,400$ a day.
\textbf{4.} $24 \times 4/7 = 13.7$ million shares a day.
\textbf{5.} From 24 to 30 million the padded shares lose 27 and earn the new
rebate of 28: each pays for itself, so the break-even volume is zero.
Padding gains $30 \times 0.0028 - 6 \times 0.0027 - 24 \times 0.0023 = +\$12\,600$ a
day.
\textbf{6.} $0.8 \times 0.0028 = \$0.00224$.
\textbf{7.} Client: $0.00224 \times 20$ million $\times 21 = \$940\,800$. Member:
$840\,000 - 25\,000 = \$815\,000$. The client route is better by \$125\,800.
\textbf{8.} Member, at its own 0.25\% tier: $0.0028 \times 30$ million $\times 21 -
25\,000 = \$1\,739\,000$. Client: \$1\,411\,200. Membership is better by
\$327\,800.
\textbf{9.} Direct connectivity on its own terms (the broker's risk checks
and ports are in the path), eligibility for liquidity-provider programmes,
control over order types and sessions, and confidentiality: the broker sees
its flow.
\textbf{10.} Registration, net capital, reporting, supervision, examinations,
the exchange's rulebook and market-access obligations of its own.
\textbf{11.} $(24 \times 21 - 23 \times 15)/6 = 26.5$ million a day.
\textbf{12.} $0.20\% \times 13$ billion $= 26$ million: the threshold has moved
up by 2 million shares a day without the firm doing anything. Targets must
be set on a forecast of consolidated volume and tracked daily.
\textbf{13.} At the marginal rate including the tier effect: a strategy whose
volume takes the firm over a threshold is worth more to the firm than its own
P\&L shows, and the backtest should report both the stand-alone figure and
the firm-level one.
\textbf{14.} Cliffs make volume sticky: a member near a threshold will
concentrate its flow on that exchange to cross it, and a member just over it
has a reason to stay. A smooth schedule would give up that pull.
\textbf{15.} It adds displayed liquidity, which is what the exchange pays
for, but liquidity that exists for the fee and not for the price is the
first to vanish; and it distorts routing statistics. Trades with no economic
purpose other than the rebate, or arranged between related parties, cross
into wash trading, which is prohibited.
\textbf{16.} $0.0031 - 0.0020 = 11$ cents per 100 shares. It can quote one tick
tighter in some stocks where the smaller firm cannot, or hold the same quote
with a fatter margin: in a business with a gross edge of 10 to 20 cents per
100 shares, that is most of it.
\textbf{17.} Rebates fund the spread that passive strategies earn; a cap
lowers both the rebate and the access fee, compresses the differences
between venue types, reduces the value of tiers and of scale, and shifts
revenue from rebate-driven market making to strategies that earn the spread
itself (see \cref{ch:m1:us-equity-market-structure} for the state of that
rule).
\textbf{18.} Shares per day by venue, added and removed; message rates;
average long and short balances and hard-to-borrow usage; the tier it would
reach alone; and its realistic alternative.
\textbf{19.} \textbf{13.7 million shares a day.}
\textbf{20.} Volume, measured the way the exchange chose to measure it.
\end{solution}

\begin{solution}{iq:m1:getting-access-equities:1}
The \omterm{def:m1:getting-access-equities:brokers}{executing broker} is the member whose name and risk controls carry the
firm's orders to the venue. The \omterm{def:m1:getting-access-equities:brokers}{clearing broker} is the clearing-house member
that settles the trades, holds the positions and cash, finances them and
calls margin. They can be the same firm; with a \omterm{def:m1:getting-access-equities:brokers}{give-up} agreement a firm
executes through several brokers and clears at one.
\iqlookfor{who bears the credit risk: the \omterm{def:m1:getting-access-equities:brokers}{clearing broker}.}
\end{solution}

\begin{solution}{iq:m1:getting-access-equities:2}
The exchange pays a rebate to the order that was resting (the maker) and
charges the order that traded against it (the taker), keeping the difference.
A \omterm{def:m1:getting-access-equities:tier}{volume tier} gives a better rebate or fee to members whose monthly volume
passes a threshold, usually a share of consolidated volume, and applies it to
the whole month.
\iqlookfor{the retroactive application, hence cliffs.}
\end{solution}

\begin{solution}{iq:m1:getting-access-equities:3}
For the queue. Where adding costs money few add, so a new order is near the
front and trades first; and takers are paid there, so fee-sensitive routers
hit that venue first. A faster, more certain fill against less informed flow
can be worth more than the fee difference.
\iqlookfor{queue priority as the thing being bought.}
\end{solution}

\begin{solution}{iq:m1:getting-access-equities:4}
Per fill, by venue, by add or remove, by order type, at the rate of the tier
the firm actually has, with the schedule as of the trade date (schedules
change monthly). Add clearing and regulatory fees. For capacity studies use
the marginal tier the strategy's own volume would reach. Never a flat
``half a cent a share''.
\iqlookfor{point-in-time schedules and tier awareness.}
\end{solution}

\begin{solution}{iq:m1:getting-access-equities:5}
Start as the client of a broker offering direct access and good pass-through
of its tier: days to set up, no registration. Choose the \omterm{def:m1:getting-access-equities:brokers}{clearing broker}
first: margin method, financing and borrow rates, stability. Add co-location
and \omterm{def:m1:us-equity-market-structure:sip}{direct feeds} when latency starts to matter; budget data licences,
including non-display. Register and take memberships when own-tier rebates,
programme eligibility and control outweigh capital, compliance and fixed
fees, typically at tens of millions of shares a day.
\iqlookfor{clearing first; membership as a volume decision.}
\end{solution}

\begin{solution}{iq:m1:getting-access-equities:6}
Compute the shortfall in shares, given the latest forecast of consolidated
volume, and the value of the tier on the whole month. Then: route more of
the firm's natural passive flow to that exchange; quote slightly more
aggressively where the expected loss per share is below the threshold of
\cref{prop:m1:getting-access-equities:padding}; bring forward flow that would
trade anyway. Not: trade with oneself or an affiliate, or enter orders with
no purpose but the count. Decide a stop: if the shortfall cannot be closed at
an acceptable loss by the 28th, stop padding.
\iqlookfor{the shadow price, and the wash-trading boundary.}
\end{solution}
